% BHU PYQ -- Transcribed & Corrected
% Paper: JYFMD21 / JYFMD-21: Jyotish (Falit)
% Exam: Under Graduate Semester II Examination, 2024-25 (NEP)
% Category: Multidisciplinary Course (MD)
% Department: Jyotish

\documentclass[11pt,a4paper]{article}
\usepackage[a4paper,top=2.5cm,bottom=2.5cm,left=2.8cm,right=2.8cm,headheight=14pt]{geometry}
\usepackage{amsmath,amssymb}
\usepackage[T1]{fontenc}
\usepackage[utf8]{inputenc}
\usepackage{lmodern,microtype,fancyhdr,enumitem,hyperref,lastpage,parskip}

\hypersetup{
    colorlinks=true,
    urlcolor=black,
    linkcolor=black,
    pdftitle={JYFMD21: Jyotish (Falit) -- BHU Sem II 2024-25 (NEP MD)}
}

\pagestyle{fancy}
\fancyhf{}
\renewcommand{\headrulewidth}{0.4pt}
\renewcommand{\footrulewidth}{0.4pt}
\fancyhead[L]{\small   \textit{Banaras Hindu University}}
\fancyhead[R]{\small   \textit{JYFMD21: Jyotish (Falit) (MD)}}
\fancyfoot[C]{\small Page  \thepage\ of \pageref{LastPage}}


\begin{document}


\begin{center}
  {\large\bfseries BANARAS HINDU UNIVERSITY}\\[2pt]
  {\normalsize Under Graduate Semester --- II Examination, 2024-25 (NEP)}\\[6pt]
  \rule{0.8\linewidth}{0.5pt}\\[6pt]
  {\large\bfseries JYOTISH (FALIT)}\\[4pt]
  {\large\bfseries Paper Code: JYFMD-21}\\[4pt]
  {\normalsize Time: 2:30 Hours\hfill Maximum Marks: 70}
\end{center}

\rule{\linewidth}{0.4pt}

\smallskip



\noindent (Write your Roll No. at the top immediately on the receipt of this question paper)

\smallskip



\noindent   \textit{Note: Answer any four questions. All questions carry equal marks.\hfill [$4  \times 17.5 = 70$ Marks]}

\bigskip


\begin{enumerate}[label=   \textbf{\arabic*.}]
    \item Give an introduction to the five components of Panchanga.

    \item Write the names of the Rasis starting from Mesha, the months starting from Chaitra, and the Nakshatras starting from Ashwini in proper sequence.

    \item Present the process of Lagna Sadhana in Kashi with a self-assumed (hypothetical) example.

    \item Define Ishtakala with an example. If the sunrise on 11/05/2025 is at 05:20 AM, calculate the Ishtakala for the following birth times:
    
\begin{enumerate}[label=(\alph*)]
        \item Birth time: 11/05/2025, 10:45 AM
        \item Birth time: 11/05/2025, 9:35 PM
        \item Birth time: 12/05/2025, 04:26 AM
    \end{enumerate}

    \item Write the method for Graha Spashta (calculating planetary positions) and present a hypothetical example for one planet.

    \item Write the method for Chandra Spashta (calculating the longitude of the Moon) with a hypothetical example.

    \item Write a general introduction to Dasha.

    \item Write the calculation method for Bhukta and Bhogya in Vimshottari Dasha with an example.
\end{enumerate}

\vfill
\rule{\linewidth}{0.4pt}

\begin{center}\small
\href{https://bhu-exam.vercel.app/}{Click here to visit our website}\\[4pt]
\href{https://me-aryan.vercel.app/}{Click here to visit our contributor}
\end{center}

\end{document}
